%==============================================================================
\chapter{Context and problem definition}
%==============================================================================

\section{General overview}
The \textbf{Zümm} project aims to design an information system dedicated to the
management and monitoring of beekeeping activities. Modern beekeeping relies on
the regular monitoring of numerous hives, often spread across several
geographic sites and operated by different beekeepers. Tracking colony health,
productivity and maintenance operations represents a significant organisational
burden, today handled in a largely manual and fragmented way.

\section{Problem statement}
Without a centralised tool, beekeeping operators face several difficulties:
\begin{itemize}
  \item no structured traceability of hive visits and inspections
  \item difficulty in planning and coordinating the work of beekeeping agents
  \item lack of visibility over production (honey, weight) and over the profitability of each site or hive
  \item late detection of anomalies (population decline, production drop, unplanned opening of a hive)
  \item inability to automatically process measurements from heterogeneous sensors.
\end{itemize}

\section{Purpose}
The system must provide a structured response in two complementary parts:
\begin{enumerate}
  \item a \textbf{manual management} of handling operations, their planning and
        monitoring (first part, the subject of this project)
  \item a \textbf{remote, automated supervision} of performance indicators,
        fed by sensors, intended to be developed as part of a later project but
        whose integration the system must anticipate.
\end{enumerate}
